\section{Previous work} \label{sec:previos} 

The problem addressed here is related to the exploitation of hydrocarbons, particularly the oil industry in Argentine Patagonia. No works have been found that address the problem exactly as described here; therefore, some works with similar aspects are listed below.

In \cite{aronofsky1962,hartsock1971}, the importance of applying optimization methods and developing operational research strategies in the oil industry is introduced, within the context of limited operating resources and during a period of significant industry growth.

In \cite{duhamel2009multicommodity}, the authors present the "Mobile Oil Recovery" problem for the Petrobras company. This involves obtaining oil by pumping from a mobile unit (truck) that extracts from a field with low-yield wells. The optimization problem consists of maximizing extraction while reducing costs. The paper proposes the problem with a single pumping unit or with several. The proposed solution was tested with up to 200 wells in the battery. The fundamental difference with the present work is that this does not involve a mobile extraction unit, but the operator must travel through the well network to adjust the production parameters of each one according to the calculated optimal value.

In \cite{langvik2012optimization}, an optimization model for oil production on offshore platforms is presented. The objective of the work is to maximize oil extraction, considering the yield in relation to the water, crude oil, and gas extracted from each well, as well as the storage and routing of the generated flows. A linearization of the model is performed using mixed-integer linear programming techniques.

This work combines the optimization of the working regimes of the well battery to obtain a net production, minimizing loss on one hand and the path the operator must travel to reconfigure the wells on the other. During the bibliographic survey recently conducted by the authors, no model proposals like the one presented in this work have been found.